% !TeX root = ../main.tex

\documentclass[../main.tex]{subfiles}
\begin{document}

\begin{center}
    \Large{\textbf{MỞ ĐẦU}}
\end{center}
\vspace{0.5cm}

Tại thành phố Hội An, lượng chất thải rắn sinh hoạt (CTRSH) dự báo đạt khoảng 120--140 tấn/ngày trong giai đoạn 2025--2030, với thành phần biến động mạnh theo mùa du lịch. Trong khi đó, chôn lấp vẫn là phương thức xử lý chủ đạo tại đây, gây áp lực lên quỹ đất hạn hẹp của thành phố ven biển, tiềm ẩn nguy cơ ô nhiễm nước ngầm và không khí tại khu vực nhạy cảm về di sản và du lịch, đồng thời bỏ qua nguồn năng lượng tiềm năng chứa trong chất thải. Chuyển sang công nghệ xử lý nhiệt có khả năng thu hồi tài nguyên và giảm mạnh tỷ lệ chôn lấp là yêu cầu cấp thiết đối với địa phương này.

Dự án Nhà máy xử lý CTRSH thành phố Hội An ban đầu được thiết kế theo công nghệ khí hóa -- tuabin khí. Tuy nhiên, qua quá trình đánh giá thực tế, công nghệ này tỏ ra kém phù hợp với nguồn rác có thành phần ẩm cao và biến động lớn của Hội An, đồng thời yêu cầu vận hành phức tạp và chi phí đầu tư cao hơn. Xuất phát từ đó, đồ án thực hiện đề tài \textit{"Nghiên cứu công nghệ xử lý rác 120 tấn/ngày đêm, ứng dụng phát điện"}, tập trung vào việc luận chứng và tính toán cho phương án chuyển đổi sang công nghệ nhiệt phân lò quay liên tục kết hợp chưng cất dầu và phát điện bằng động cơ diesel -- giải pháp linh hoạt hơn với biên độ vận hành rộng và sản phẩm đầu ra đa dạng.

Các mục tiêu nghiên cứu cụ thể bao gồm: (i) lựa chọn và luận chứng phương án công nghệ phù hợp với điều kiện địa phương; (ii) tính toán cân bằng vật chất và cân bằng năng lượng cho hệ thống công suất 120 tấn/ngày đêm; (iii) xác định cấu hình và thông số cơ bản của các thiết bị chính; (iv) đánh giá hiệu quả kinh tế thông qua CAPEX, OPEX, NPV và IRR; (v) đề xuất giải pháp kiểm soát môi trường theo QCVN 19:2024/BTNMT, QCVN 40:2025/BTNMT và bảo đảm an toàn vận hành theo quy định hiện hành.

Phạm vi nghiên cứu được giới hạn trong hệ thống nhà máy xử lý CTRSH công suất thiết kế 120 tấn/ngày đêm tại thành phố Hội An. Các thông số đầu vào như thành phần rác, độ ẩm và công suất được lấy từ tài liệu dự án; các tham số kỹ thuật về tỷ lệ phân bố sản phẩm nhiệt phân và hiệu suất chuyển hóa được tham khảo từ nghiên cứu quốc tế về nhiệt phân chất thải rắn đô thị.

Đồ án được tổ chức thành năm chương. Chương 1 trình bày tổng quan về chất thải rắn sinh hoạt, đặc điểm rác thải tại Hội An và các công nghệ xử lý nhiệt. Chương 2 xây dựng cơ sở lý thuyết về nhiệt phân, chưng cất dầu và phát điện diesel, đồng thời luận chứng lựa chọn công nghệ. Chương 3 thực hiện các tính toán cân bằng vật chất, cân bằng năng lượng và thiết kế sơ bộ thiết bị chính. Chương 4 đánh giá các hệ thống phụ trợ, hiệu quả kinh tế, an toàn công nghiệp và kiểm soát môi trường. Chương 5 tổng hợp kết quả nghiên cứu, đánh giá tính khả thi của phương án và đề xuất hướng phát triển tiếp theo.

\end{document}